\section{Benchmark Protocol}
\label{sec:bench}

\subsection{Splits and Generalization Controls}
\label{sec:bench-splits}

The benchmark publishes a canonical split of benign and anomalous trajectories per environment, stratified so that no task template appears in both train and test. Beyond the canonical split, the protocol defines three controls that respond directly to documented failure modes of prior safety evaluation. First, \emph{leave-one-attack-family-out} (LOAO): for each family $a \in \{A1, A2, A3\}$, detectors are trained on benign data plus the two other families and evaluated on $a$, measuring whether a detector tuned on known threats generalizes to unseen ones---the property that~\cite{liu2026passing} shows prompt-injection detectors lose when evaluation distributions shift. Second, \emph{cross-environment transfer}: detectors are trained on replayed environments and evaluated on a held-out environment whose tool catalog and task composition differ, with no re-tuning. Third, \emph{contamination control}: the canonical split is constructed so that textual content of benign and anomalous trajectories is paraphrase-balanced, so that a text classifier cannot separate families by surface style alone; we verify this by reporting a text-only logistic-regression probe on $D_1$ inputs.

\subsection{Mimicry Variants}
\label{sec:bench-mimicry}

Provenance-graph detectors in systems security are defeated by mimicry: adversaries that preserve benign structural statistics while executing the malicious payload~\cite{goyal2023sometimes}. Our protocol instantiates three mimicry variants over A1 and A3 trajectories. \emph{M1 (structure splicing)} interleaves benign calls from matched benign trajectories around the malicious calls, preserving degree profiles and temporal regularity. \emph{M2 (padding)} appends benign tail calls after the payload so that whole-graph statistics (counts, tool histograms) regress toward the benign distribution. \emph{M3 (argument camouflage)} rewrites argument entities---splitting exfiltration across chunks, renaming accounts to benign-looking strings---so that entity-level features lose their signal while the dataflow graph is preserved. Because M3 changes entities but not edges, the comparison between M2 and M3 degradation cleanly attributes failure to surface features rather than structure. All variants are generated by deterministic scripts and released with the benchmark, so reported mimicry results are reproducible transformations, not researcher degrees of freedom.

\subsection{Metrics}
\label{sec:bench-metrics}

Detection is scored with AUROC and AUPRC (the latter primary, given class imbalance), plus Recall@$k$ at the operator-relevant operating points $k \in \{1\%, 5\%\}$ of the ranked audit queue and FPR@95TPR, which measures the false-alarm burden that a deployable guardrail imposes. Localization on A1 is scored as node-level AUC against step labels~\cite{pinjari2026agentdrift}. Generalization is summarized by the LOAO drop---the AUPRC gap between seen-family and unseen-family evaluation---and transfer is scored as the retained fraction of in-distribution AUPRC on the held-out environment. Audit economics reports GPU-seconds per trace and, where applicable, API cost per trace for LLM-judge baselines, because a detector that is accurate but costs more than the incident it prevents is not an auditing solution. All results are reported as mean $\pm$ standard deviation over five seeds; family-level comparisons use two-sided Wilcoxon signed-rank tests over seeds with Holm correction across the detector matrix.

\subsection{Fairness Rules}
\label{sec:bench-fairness}

Three rules keep the comparison meaningful. First, \emph{identical inputs}: every detector consumes the same compiled graphs at the same disclosure level; no baseline may read the raw trace. Second, \emph{tuning budget parity}: hyperparameter search is bounded identically across detectors, and the search space is published. Third, \emph{no family peeking}: unsupervised methods never see family labels; supervised methods see them only under the LOAO restriction. The reference detector is held to the same rules as every baseline---its role is to be a strong, reproducible point on the detector map, not an advertised winner.
